\documentclass[12pt]{article}
\usepackage[margin=1in]{geometry} % set suitable margins
\usepackage{amsmath}
\usepackage{amssymb}
\numberwithin{equation}{section}
\title{EMCH 354 Extra Exercise}
\author{J. Vaught}
\date{}

\begin{document}
\maketitle

\section{Problem \#1}
To derive the thermal resistance of cylindrical shell structures from the thermal diffusion equation in the Cylindrical Coordinate System, we start with the general heat conduction equation in cylindrical coordinates without internal heat generation:
\begin{equation}
    \frac{1}{r} \frac{\partial}{\partial r}\left( k r \frac{\partial T}{\partial r} \right) + \frac{1}{r^2} \frac{\partial}{\partial \theta}\left( k \frac{\partial T}{\partial \theta} \right) + \frac{\partial}{\partial z}\left( k \frac{\partial T}{\partial z} \right) = \rho c_p \frac{\partial T}{\partial t}
\end{equation}

Where:
\begin{itemize}
    \item \( T \) is the temperature,
    \item \( r \) is the radial coordinate,
    \item \( \theta \) is the angular coordinate,
    \item \( z \) is the axial coordinate,
    \item \( k \) is the thermal conductivity,
    \item \( \rho \) is the density,
    \item \( c_p \) is the specific heat capacity,
    \item \( t \) is time.
\end{itemize}

For a steady-state condition and assuming that the temperature gradient in the angular (\( \theta \)) and axial (\( z \)) directions are negligible, the equation simplifies to:
\begin{equation}
    \frac{1}{r} \frac{\partial}{\partial r}\left( k r \frac{\partial T}{\partial r} \right) = 0
\end{equation}

Integrating this equation with respect to \( r \), we get:
\begin{equation}
    k r \frac{\partial T}{\partial r} = C_1
\end{equation}

Where \( C_1 \) is the integration constant. To find \( C_1 \), we apply the boundary conditions. If we consider a cylindrical shell with inner radius \( r_i \) and outer radius \( r_o \), and temperatures \( T_i \) and \( T_o \) respectively, we can integrate again to find the temperature distribution:
\begin{equation}
    T_o - T_i = \frac{C_1}{k} \ln\left(\frac{r_o}{r_i}\right)
\end{equation}

The thermal resistance \( R_{th} \) for conduction through a cylindrical shell is defined as the temperature difference divided by the heat transfer rate (\( Q \)):
\begin{equation}
    R_{th} = \frac{\ln\left(\frac{r_o}{r_i}\right)}{2 \pi k L}
\end{equation}

Where \( L \) is the length of the cylindrical shell. This is the expression for the thermal resistance of a cylindrical shell structure in the radial direction for steady-state heat conduction.
\newline\newline
**Please note that this derivation assumes one-dimensional heat flow and neglects any heat transfer in the angular and axial directions. For more complex situations involving multiple layers or different boundary conditions, the analysis would need to be expanded accordingly.

\newpage
\section{Problem \#2}

Given a spherical, stainless steel canister with reacting chemicals providing a uniform heat flux $q''_i$ to its inner surface, we are tasked with finding the initial rate of change of the wall temperature and the steady-state temperature.

\subsection*{Given Data}
The given data for the stainless steel canister are:
\begin{align*}
    r_o & = 0.6 \, \text{m} \\
    r_i & = 0.5 \, \text{m} \\
    T_i & = 500 \, \text{K} \\
    \rho & = 8055 \, \text{kg/m}^3 \\
    c_p & = 510 \, \text{J/kg} \cdot \text{K}
\end{align*}


\subsection*{Objective}
The objective is to determine:
\begin{enumerate}
    \item The initial rate of change of the wall temperature.
    \item The steady-state temperature of the wall.
\end{enumerate}

\subsection*{Properties of Stainless Steel (AISI 302)}
Stainless steel (AISI 302) has a density $\rho$ of $8055 \, \text{kg/m}^3$ and a specific heat capacity $c_p$ of $510 \, \text{J/(kg \cdot K)}$.

\subsection*{Assumptions}
The following assumptions are made in the heat transfer analysis:
\begin{itemize}
    \item The temperature gradient in the wall is neglected.
    \item The properties of the material are constant.
    \item The time-dependent heat flux is uniform at the inner surface.
\end{itemize}

\subsection*{Part (a) Initial Rate of Change of Wall Temperature}
The heat transfer equation, neglecting the work done and with constant properties, is given by the first law of thermodynamics:
\begin{equation}
    mc_p \frac{dT}{dt} = q''_i A_i - h A_o(T - T_{\infty})
\end{equation}
where:
\begin{align*}
    A_i &= 4\pi r_i^2 \\
    A_o &= 4\pi r_o^2 \\
\end{align*}
The mass $m$ of the canister wall is:
\begin{equation}
    m = \frac{4}{3}\pi (r_o^3 - r_i^3) \rho
\end{equation}
Substituting $A_i$, $A_o$, and $m$ into the heat transfer equation, we get:
\begin{equation}
    \frac{4}{3}\pi (r_o^3 - r_i^3) \rho c_p \frac{dT}{dt} = q''_i 4\pi r_i^2 - h 4\pi r_o^2 (T_i - T_{\infty})
\end{equation}
Solving for the rate of change of temperature at the initial moment, we find:
\begin{equation}
    \left.\frac{dT}{dt}\right|_{t=0} = -0.08827 \, \text{K/s}
\end{equation}

\subsection*{Part (b) Steady-State Temperature of the Wall}
For the steady-state condition, the rate of change of temperature with time is zero ($\frac{dT}{dt} = 0$), and we have:
\begin{equation}
    q''_i A_i = h A_o (T_{ss} - T_{\infty})
\end{equation}
where $T_{ss}$ is the steady-state temperature. Solving for $T_{ss}$ yields:
\begin{equation}
    T_{ss} = T_{\infty} + \frac{q''_i r_i^2}{h r_o^2}
\end{equation}
By substituting the known values, we can calculate the steady-state temperature as follows:
\begin{equation}
    T_{ss} = 300 \, \text{K} + \frac{10^5 \times 0.5^2}{500 \times 0.6^2}
\end{equation}
This simplifies to:
\begin{equation}
    T_{ss} = 438.889 \, \text{K}
\end{equation}



\end{document}
